problem:  
Let \( X \) be a Calabi–Yau threefold defined over \( \mathbb{C} \), and let \( \mathcal{M}_{\mathrm{complex}}(X) \) denote its complex structure moduli space equipped with the Weil–Petersson metric induced by variations of Hodge structure. Let \( X^\vee \) denote the mirror manifold, and let \( \mathcal{M}_{\mathrm{Kähler}}(X^\vee) \) be its Kähler moduli space endowed with the quantum-corrected special Kähler metric defined via the genus-zero Gromov–Witten potential.  

It is known that near a large complex structure limit point, the mirror map identifies the asymptotic special Kähler structures on neighborhoods of the boundary in \( \mathcal{M}_{\mathrm{complex}}(X) \) and \( \mathcal{M}_{\mathrm{Kähler}}(X^\vee) \). However, globally the mirror map is neither known to be an isometry nor to extend continuously across the compactified moduli spaces due to the presence of singularities and instanton corrections.  

**Problem:**  
Is there a canonical metric completion—possibly in the sense of a hybrid or non‑archimedean analytic space à la Kontsevich–Soibelman or of the tropical compactification in the Gross–Siebert program—of the special Kähler metrics on both \( \mathcal{M}_{\mathrm{complex}}(X) \) and \( \mathcal{M}_{\mathrm{Kähler}}(X^\vee) \) such that the mirror map extends to a global isometry between these completed spaces?  
Equivalently, can one construct a “mirror symmetric” metric completion \( \overline{\mathcal{M}}_{\mathrm{mirror}} \) equipped with two canonical projections to the complex and quantum-corrected Kähler moduli spaces that are local isometries away from the boundary and extend continuously to isometric identifications on the boundary strata determined by limiting mixed Hodge structures?  

Why is it a "good" problem:  
This refined problem retains mathematical precision while addressing the fundamental obstruction to realizing mirror symmetry as a *global metric equivalence*. It unifies several major frameworks—classical variation of Hodge structures, Gromov–Witten theory, degeneration theory of Calabi–Yau manifolds, and non‑archimedean or tropical compactifications—into a single question about the existence of a canonical metric completion compatible with mirror duality. The problem is research-level and unresolved: none of the existing approaches (SYZ fibrations, Gross–Siebert reconstruction, or non‑archimedean mirror symmetry) currently yield such a global isometric correspondence. Its conclusion, if true, would integrate analytic, geometric, and arithmetic aspects of mirror symmetry, providing a new geometric object encoding the full “completed” equivalence between complex and Kähler moduli spaces. The formulation is simple yet probes deeply into the unification of differential geometry, Hodge theory, and enumerative geometry.